\documentclass[11pt]{article}
\input{../template/preamble}
\usepackage{longtable,array,tabularx}

\title{Estimer le besoin de soins quand on n'observe que le recours :\\
des estimateurs concurrents, sept protocoles gelés, et un indice de couverture qui sait dire ce qu'il ignore}
\author[1]{Équipe hackathon-santé}
\affil[1]{Projet « jumeau numérique de l'offre de soins », dépôt \texttt{hackathon-sante/v2}}
\date{Version du \today}

\begin{document}
\sloppy
\maketitle

\begin{abstract}
\textbf{Contexte.} Le jumeau numérique projette jusqu'en 2050 l'effectif du secteur privé de cinq professions de santé
(sages-femmes, chirurgiens-dentistes, pharmaciens, masseurs-kinésithérapeutes, infirmiers), de la commune au national.
Pour dire si cet effectif suffira, il faut un besoin. Or le besoin n'est jamais observé : on observe le recours (actes
remboursés), qui dépend de l'offre.

\textbf{Méthode.} Nous formalisons cette censure : $R = B\,\varphi(A) + D(A)$, où $B$ est le besoin, $\varphi$ la part
satisfaite selon l'accès $A$, et $D$ la demande induite. Sans restriction, $B$ n'est pas identifiable. Six
estimateurs, chacun porteur d'une restriction déclarée, partagent une même interface et un même indice de couverture
indexé à 2026. Leur survie se décide par sept protocoles écrits et gelés avant tout calcul, avec une empreinte
horodatée par deux autorités RFC~3161.

\textbf{Résultats.}
\begin{itemize}
  \item La pondération par l'âge (M1) est retenue pour les cinq professions. Pour les infirmiers, la part des
        départements où elle déplace l'indice 2050 de moins de 2~\% vaut 0~\%\mesure{}.
  \item La saturation du recours (M2b) est rejetée pour les quatre professions testées : le besoin non satisfait
        n'est pas identifié en données ouvertes.
  \item La morbidité (M3) ne modifie l'indice que pour les infirmiers.
  \item Au backtest 2016--2025, l'âge bat la persistance pour trois professions sur quatre (infirmiers :
        $-8{,}1$\mesure{} points de MAPE).
  \item L'accessibilité spatiale (M4, famille 2SFCA) ne reproduit pas l'APL de la DREES (Spearman 0,71 à
        0,76\mesure{}, même avec des temps routiers), mais bat nettement la densité simple.
  \item Pour le transfert « vaccination en officine », la montée en charge est mesurée
        ($s_\infty = 0{,}74$\derive{}, $\tau = 2{,}4$\derive{} ans). Aucun effet sur le total des vaccinations n'est
        détectable.
\end{itemize}

\textbf{Limite principale.} L'indice mesure la couverture de la \emph{demande adressée au secteur privé de ville},
pas le besoin de toute la population. En 2050, le choix de l'estimateur pèse plus que la politique simulée pour les
kinés et les infirmiers : leur indice s'affiche en fourchette.
\end{abstract}
\noindent\textbf{Mots-clés :} besoin de soins ; recours ; demande induite ; identification ; 2SFCA ; APL ; transfert
de tâches ; pré-enregistrement.

\noindent\small Exposants de provenance : \mesure{} mesuré, \derive{} dérivé, \suppose{} supposé. Les valeurs
sourcées portent leur référence.\normalsize

% ---------------------------------------------------------------------------------------------------
\section{Introduction}
La question type du sujet est : \enquote{combien de spécialistes à Clermont-Ferrand en 2035 si les pharmaciens peuvent
vacciner ?}. Elle contient deux questions. \emph{Combien} : le jumeau y répond par ses projections d'effectifs.
\emph{Est-ce assez} : il faut un dénominateur, le besoin. Ce dénominateur rend aussi calculables les transferts de
tâches, puisqu'un transfert ne crée aucun professionnel : il déplace une part du besoin d'une profession vers une
autre.

Jusqu'ici, la carte affichait une densité pour \num{100000} habitants. Cela revient à supposer un besoin identique par
habitant, ce qui ignore le vieillissement et rend les transferts invisibles.

\paragraph{Contributions.}
\begin{enumerate}
  \item Une formalisation du lien recours/besoin, et la preuve que le besoin n'est pas identifié sans restriction.
  \item Six estimateurs concurrents derrière une interface commune, chacun avec sa restriction, son cas limite et son
        test d'abandon.
  \item Sept protocoles pré-enregistrés, dont l'antériorité est prouvée hors de git par horodatage externe.
  \item Un opérateur de transfert de tâches conservatif, qui sépare la conversion de temps $k$ de l'expansion du
        recours $m$, calé sur le déploiement échelonné de la vaccination en officine.
  \item À notre connaissance, et sous réserve de la recherche journalisée, une projection d'un indice de type APL à
        2050 sous scénarios conjoints de démographie et d'effectifs, avec un écart mesuré à l'APL publiée.
\end{enumerate}

% ---------------------------------------------------------------------------------------------------
\section{Travaux existants}\label{sec:litterature}
Une passe de littérature a été faite par trois lecteurs, le plus souvent sur résumés, avec les DOI vérifiés. Une
recherche de nouveauté journalisée (14 requêtes) l'a complétée. Rien ici n'est revendiqué comme \enquote{premier}.

\paragraph{Demande induite.} L'idée que l'offre crée une partie de son recours remonte à \citet{evans1974}. En France,
une étude en panel sur \num{4500} médecins libéraux de 1979 à 1993 trouve qu'une densité plus forte baisse à peine
l'activité par médecin, compensée par le volume par séance~\citep{delattre2003}. Une revue de 25~études associe
toujours densité et consommation, sans saturation rapportée~\citep{leonard2009}. C'est l'a priori défavorable à M2b,
confirmé par P-B2.

\paragraph{Neutraliser l'offre.} Les formules d'allocation des budgets de santé en Angleterre régressent le recours sur
des indicateurs de besoin \emph{et} d'offre, puis l'évaluent à l'offre moyenne nationale~\citep{carrhill1994,
gravelle2003}. M2a reprend cette méthode. Remplacer la référence \enquote{offre moyenne} par un niveau de saturation
estimé, avec une règle d'abandon en l'absence de plateau, est l'extension testée ici (M2b) ; nous ne l'avons pas
trouvée publiée.

\paragraph{Besoin épidémiologique.} Le cadre de \citet{birch2007} écrit l'effectif requis comme le produit de la
population, de l'état de santé, du niveau de service par besoin et de la productivité.
\citet{mackenzie2019} l'étendent à plusieurs professions, avec des parts de services et des productivités par
profession. M3 applique ce cadre à la cartographie des pathologies de la CNAM. L'opérateur de transfert en est proche,
mais il ajoute l'expansion $m$ et une montée en charge $s(t)$.

\paragraph{Accessibilité spatiale.} Le 2SFCA de \citet{luowang2003} et sa version à paliers~\citep{luo2009} fondent
l'APL de la DREES et de l'IRDES~\citep{barlet2012,vergier2017}. Les paliers de l'APL sont une convention négociée, pas
une estimation~\citep{vergier2017}. Les formes \enquote{équilibrées} préservent les totaux~\citep{paez2019}. Les
travaux prospectifs trouvés projettent la population à offre figée. Nous n'avons trouvé ni projection de l'APL sous
scénarios conjoints de démographie et d'effectifs, ni reproduction publiée de l'APL en données ouvertes.

\paragraph{Transferts de tâches.} La substitution médecin $\to$ infirmier allonge probablement les consultations
($k>1$) et augmente un peu les retours ($m>0$)~\citep{laurant2018}. Au Canada, l'autorisation de vacciner en officine
n'augmente la couverture que de 5~\% en valeur relative~\citep{buchan2017}. En France, la part des pharmaciens en
2019--20 suit la date d'entrée des régions dans l'expérimentation~\citep{liard2022}.

% ---------------------------------------------------------------------------------------------------
\section{Le problème : un besoin censuré par l'offre}\label{sec:probleme}
Pour une cellule $i$ (territoire, âge, sexe) et une profession $p$ :
\[
R_{ip} = B_{ip}\,\varphi_p(A_{ip}) + D_p(A_{ip}) + \varepsilon_{ip},\qquad 0<\varphi\le 1,\quad D\ge 0 .
\]
\begin{itemize}
  \item $R$ est le recours observé, $B$ le besoin latent selon une norme déclarée, $A$ l'accessibilité de l'offre.
  \item $\varphi$ est la part du besoin satisfaite. Elle reflète le rationnement par la distance, les délais ou le
        renoncement.
  \item $D$ est la demande induite par l'offre~\citep{evans1974}. En France, une étude sur des médecins libéraux de
        1979 à 1993 en trouve un signe fort~\citep{delattre2003} ; une revue de 25~études associe toujours densité et
        consommation~\citep{leonard2009}.
\end{itemize}

\paragraph{Non-identification.} Avec $D=0$, pour toute fonction $h>0$, le couple $(B\,h,\ \varphi/h)$ produit le même
$R$\derive{}. Le niveau du besoin et la forme du rationnement ne sont donc pas séparables à partir de $(R,A)$. Il faut
une restriction parmi trois :
\begin{itemize}
  \item l'indépendance conditionnelle, $B\perp A\mid X$ ;
  \item la saturation, $\varphi(A)\to1$ quand $A\to\infty$ ;
  \item une exclusion, c'est-à-dire une variable qui déplace l'accès sans toucher au besoin.
\end{itemize}
Une quatrième voie remplace l'identification par une norme externe à l'offre.

\paragraph{Ce que l'indexation sauve.} L'indice est
\[
C(p,T,t)=\frac{\text{offre}(p,T,t)\,/\,B(p,T,t)}{\text{offre}(p,T,2026)\,/\,B(p,T,2026)} .
\]
Tout biais multiplicatif constant par territoire s'y annule\derive{}. Seul compte le biais qui change avec le temps. On
peut donc comparer 2035 à 2026 dans un même territoire avec un estimateur simple ; comparer deux territoires entre eux
en niveau demande davantage.

\paragraph{Ce que l'indice mesure vraiment.} Le recours ouvert est celui des soins de ville remboursés. Il ne compte
pas l'hôpital public. Les estimateurs M1 à M4 mesurent donc la \emph{demande adressée au secteur privé de ville}, et
non le besoin de toute la population. L'interface et la sortie de l'API le disent (attribut \texttt{grandeur}).

% ---------------------------------------------------------------------------------------------------
\section{Estimateurs}\label{sec:estimateurs}
Tous passent par la même interface (\texttt{v2/besoin/interface.py} : \texttt{fit}, \texttt{predict},
\texttt{assumptions}, \texttt{norme}, \texttt{grandeur}). Ils produisent $B$ en équivalents temps plein, puis $C$.

\begin{table}[h]\centering\small
\begin{tabularx}{\linewidth}{@{}l X X@{}}
\toprule
& \textbf{Forme} & \textbf{Restriction qui l'identifie} \\
\midrule
M0 & $B_0=\bar r_p\,\mathrm{pop}(T,t)$ & besoin par habitant uniforme (cas limite de tous les autres) \\
M1 & $B_1=\sum_{a,s}\mathrm{pop}(T,t,a,s)\,r(p,a,s)$ & besoin par âge et sexe = recours national par âge et sexe \\
M2a & recours régressé sur besoin et offre, offre fixée à sa moyenne nationale & neutralisation de l'offre, méthode des formules de capitation anglaises~\citep{carrhill1994,gravelle2003} \\
M2b & $R_{\mathrm{std}}=\rho_{\max}A/(A+K)$, besoin au plateau & saturation (extension proposée ici) \\
M3 & $\sum_g \mathrm{pop}\cdot\mathrm{prev}(g)\cdot n(p,g)$ & norme : prévalence $\times$ soins par patient~\citep{birch2007,mackenzie2019} \\
M4 & 2SFCA : $R_j=S_j/\sum_k w(t_{kj})P_k$, $A_i=\sum_j w(t_{ij})R_j$ & aucune : indice d'accès qui prend le besoin de M1 à M3 en entrée~\citep{luowang2003,luo2009} \\
\bottomrule
\end{tabularx}
\caption{Les estimateurs du besoin.}
\end{table}

\paragraph{M1.} La \emph{forme} du profil vient des poids d'âge publiés avec l'APL de la DREES : consommation DCIR
2020--2022, 19~tranches quinquennales. Pour les infirmiers, 90~ans et plus pèsent 15,5 fois la moyenne, et 80 à
84~ans 4,5 fois. Le \emph{niveau}, et le profil des pharmaciens, viennent d'Open DAMIR 2025 (12~mois). M1 garde tous
les statuts juridiques du secteur privé de ville. Le statut inconnu représente 3,3 à 12,6~\%\mesure{} des actes selon
la profession, et l'exclure ne déplacerait la répartition par âge que de 0,2 à 2,7~\%\mesure{}.

\paragraph{M2a et M2b.} La neutralisation de l'offre est établie dans les formules d'allocation, qui évaluent le
recours à l'offre moyenne nationale~\citep{gravelle2003}. Nous proposons d'évaluer plutôt au \emph{point de
saturation}, avec une règle d'abandon : s'il n'y a pas de plateau, la demande induite domine et M2b tombe. La
littérature française laisse prévoir ce résultat~\citep{delattre2003}.

\paragraph{M3.} Les prévalences viennent de la cartographie des pathologies de la CNAM : département, âge quinquennal,
sexe, 61~groupes, 2015--2024. Les soins par patient viennent des dépenses nationales par poste et par pathologie,
converties en actes par le prix moyen d'Open DAMIR. Cette combinaison est la nôtre.

\paragraph{M4.} Les paliers et les bornes sont ceux publiés par la DREES pour l'APL 2024~\citep{vergier2017} : poids
1, 2/3, 1/3 sur 0--10, 10--15, 15--20~min (30~min pour les sages-femmes). L'offre $S_j$ est l'effectif de 65~ans au
plus, multiplié par un ETP approché, car l'activité individuelle n'est pas ouverte. Les temps de trajet viennent
d'OSRM sur OpenStreetMap, en voiture, de centre à centre. Les poids d'âge sont \emph{figés} à une année de base. La
DREES les renormalise à chaque millésime, si bien que son APL ne capte pas le vieillissement national ; figés, ils
rendent M4 projetable. La conservation $\sum_i P_iA_i=\sum_j S_j$ tient à $10^{-16}$\mesure{} près~\citep{paez2019}.

% ---------------------------------------------------------------------------------------------------
\section{Transferts de tâches}\label{sec:transferts}
Un acte passe de la profession A à B à partir de la date d'autorisation $t_0$ :
\[
B'_A=B_A(1-s(t)\,a),\qquad B'_B=B_B+s(t)\,a\,k\,B_A+m(t),\qquad \Sigma B'-\Sigma B=s\,a\,B_A(k-1)+m .
\]
\begin{itemize}
  \item $a$ est la part du temps de A consacrée à l'acte.
  \item $k=t_B/t_A$ est le rapport des durées de l'acte chez B et chez A.
  \item $s(t)=s_\infty(1-e^{-(t-t_0)/\tau})$ est la part effectivement transférée. Le levier politique est la date
        d'autorisation $t_0$, pas $s$.
  \item $m$ est l'expansion : des patients qui n'auraient pas eu l'acte du tout.
\end{itemize}
Les effectifs ne changent pas : un test du projet a mesuré un effet nul de la densité sur les installations.

La médecine générale et les spécialistes entrent comme professions \emph{source} seulement. On rapporte les actes
libérés, jamais un indice de couverture pour elles. L'identité de conservation est vérifiée à $5{,}7\times10^{-10}$
ETP\mesure{} sur \num{13947} ETP déplacés.

Les modèles de besoin à plusieurs professions répartissent déjà les services par parts et productivités, ce qui
contient un $k$ implicite~\citep{mackenzie2019}. Nous n'y avons trouvé ni $m$ ni montée en charge $s(t)$.

\paragraph{Ce qui est mesurable en données ouvertes.} Sur 25~paramètres de transfert et de besoin, 8 sont sourcés, 2
le sont en partie, 1 est supposé et 14 sont non mesurables\mesure{}. Seule la vaccination antigrippale en officine
(L23) a $a$ et $s(t)$ mesurés :
\begin{itemize}
  \item la part des pharmaciens dans les vaccinations antigrippales passe d'un peu moins d'un tiers en 2019--20 à
        près de deux tiers en 2024--25~\citep{ccomptes2026} ;
  \item $k$ reste supposé, dans la plage $[0{,}25\,;\,1]$\suppose{}, faute de durée publiée pour l'acte officinal.
\end{itemize}

\begin{figure}[h]\centering
\includegraphics[width=0.62\linewidth]{figures/s_t_vaccination.pdf}
\caption{Part des vaccinations antigrippales faites en officine : actes IQVIA de la campagne rapportés aux doses
J07BB délivrées (Open Medic) l'année de début\mesure{}. Ajustement $s(t)=s_\infty(1-e^{-(t-2019)/\tau})$ sur 2021--2024 :
$s_\infty=0{,}74\pm0{,}07$\derive{}, $\tau=2{,}4\pm0{,}5$\derive{} ans. Les fiches de l'axe~1, dérivées du SNDS,
donnaient 0,70 et 2,8~ans.}
\label{fig:st}
\end{figure}

La part officinale atteignait déjà 0,687 en 2025--26\mesure{}. Seule la part au-delà de celle atteinte en 2026
agit donc : L23 ne déplace l'indice national 2050 des pharmaciens que de 0,03~\%\mesure{} (1,0126 contre 1,0123), et
n'enlève qu'environ \num{104000} vaccinations aux généralistes en 2050\mesure{}. Une première version, qui rejouait la
montée en charge depuis 2019, en surestimait l'effet d'un facteur 19 (erratum au registre, 2026-11-15).

% ---------------------------------------------------------------------------------------------------
\section{Données}\label{sec:donnees}
\begin{table}[h]\centering\small
\begin{tabularx}{\linewidth}{@{}l X l l@{}}
\toprule
\textbf{Jeu} & \textbf{Contenu, millésime} & \textbf{Grain} & \textbf{Licence} \\
\midrule
Open DAMIR & actes et montants par profession, âge, sexe ; 12~mois de 2025 & 13 zones & LO 2.0 \\
Patientèle (data.ameli) & file active moyenne par professionnel, 2016--2025 & département & ODbL \\
Pathologies (data.ameli) & prévalence, 61 groupes, 2015--2024 ; dépenses par poste & département, national & ODbL \\
APL DREES (jeu 530) & 2022--2024 ; V0 2010/2013 pour les pharmacies & commune & LO 2.0 \\
Open Medic & vaccins antigrippaux J07BB, 2014--2025 & région & LO 2.0 \\
IQVIA (data.gouv) & actes de vaccination en officine, 2021--2026 & national & LO 2.0 \\
OpenStreetMap + OSRM & temps de trajet entre centres de communes & commune & ODbL \\
Omphale 2022 (Insee) & population par âge et sexe jusqu'en 2070 & département & LO 2.0 \\
\bottomrule
\end{tabularx}
\caption{Sources ouvertes. Aucune ne donne d'actes par profession au département.}
\end{table}

\begin{itemize}
  \item \textbf{Contrôles :} sur 1\,276 contrôles de totaux, 1\,247 passent\mesure{}. Les 29 échecs viennent du
        secret statistique et des bénéficiaires non localisés.
  \item \textbf{Open DAMIR} n'a pas d'API de requête : il est lu en flux, mois par mois (environ 0,95~Go compressé
        par mois\mesure{}), sans stocker de fichier brut.
  \item \textbf{Tables commitées :} les tables dérivées de data.ameli sont commitées sous ODbL. La matrice OSRM est
        reconstruite par script.
\end{itemize}

% ---------------------------------------------------------------------------------------------------
\section{Validation}\label{sec:validation}
\paragraph{Pré-enregistrement hors de git.} Le rapport technique a montré que l'historique git ne prouve pas
l'antériorité d'un gel. Chaque protocole est donc copié, haché (SHA-256), et l'empreinte du manifeste est horodatée
par deux autorités RFC~3161 (FreeTSA, DigiCert), puis soumise à OpenTimestamps. Seule l'empreinte sort de la machine.
Tout écart au protocole est écrit et commité \emph{avant} le calcul. Aucun seuil n'a été changé après un résultat.

\begin{longtable}{@{}p{1.3cm}p{5.2cm}p{7.6cm}@{}}
\toprule
\textbf{Test} & \textbf{Question} & \textbf{Résultat} \\
\midrule
\endhead
P-B1 & L'âge (M1) déplace-t-il l'indice 2050 par rapport à M0 ? Abandon si $S_1>90$~\% des départements sous 2~\% d'écart et $\rho\ge0{,}98$ & M1 retenu pour les 5 professions. IDE $S_1=0$~\%\mesure{}, $\rho=0{,}835$\mesure{} ; MK $S_1=0$~\%\mesure{} ; SF 1~\%\mesure{} ; PH 0~\%\mesure{} ; dentistes de justesse ($S_1=87{,}1$~\%\mesure{} avec le profil DAMIR) \\
P-B2 & Le recours standardisé plafonne-t-il quand l'accès augmente ? & M2b rejeté pour IDE, MK, CD, SF : borne haute de l'IC à 98,75~\% de $\delta\ge0$ partout\mesure{} \\
P-B3 & La morbidité (M3) apporte-t-elle au-delà de l'âge ? & Indice : M3 retenu pour les seuls IDE ($S_3=83$~\%\mesure{}). Explication du recours hors échantillon : $\Delta R^2$ de $-0{,}03$ à $+0{,}01$\mesure{} contre 0,10 exigé, rejetée \\
P-B4 & M4 reproduit-il l'APL 2024 (IDE, MK, SF, CD) ? Spearman $\ge0{,}90$, erreur médiane $\le15$~\% & Échec. Vol d'oiseau : $\rho$ 0,66 à 0,77\mesure{}. OSRM : $\rho$ 0,71 à 0,76\mesure{}, erreur médiane 23 à 33~\%\mesure{}. Gain sur la meilleure densité naïve $+0{,}15$ à $+0{,}35$\mesure{} ; accord départemental $\rho=0{,}98$\mesure{} (MK) \\
P-B4-PH & M4 reproduit-il l'APL V0 2013 des pharmacies ? & Échec : $\rho=0{,}72$ [0,71 ; 0,74]\mesure{}, erreur médiane 33~\%\mesure{}, gain $+0{,}14$\mesure{} ; niveau moyen reproduit (22,1 contre 22,5 officines pour \num{100000} hab.\mesure{}) \\
P-B5 & Backtest 2016 $\to$ 2025 : l'âge bat-il la persistance ? & Oui pour 3 sur 4 : $\Delta$MAPE IDE $-8{,}07$, MK $-3{,}32$, CD $-0{,}99$ ; SF $+1{,}57$ point\mesure{}. Le GBM concurrent n'est servi que comme prédicteur du recours \\
P-B6 & La vaccination en officine a-t-elle ajouté des vaccinations ($m>0$) ? DiD à adoption échelonnée, 4 régions traitées, 12 zones & ATT $+0{,}73$~\% [$-0{,}12$ ; $+1{,}58$]\mesure{}, permutation $p=0{,}15$\mesure{}. Placebo 0--19~ans en échec : non attribuable, $m$ reste supposé à 0 \\
\bottomrule
\caption{Les sept protocoles pré-enregistrés et leurs verdicts.}
\end{longtable}

\paragraph{Lecture de P-B4.} Le même profil revient pour toutes les professions, pharmacies comprises, où l'offre est
un simple compte d'officines : le niveau est bon, le classement fin ne l'est pas. Pour les kinés, l'accord au
département est très bon ($\rho=0{,}98$\mesure{}). L'écart se situe donc au niveau communal, dans la répartition de l'offre. Il ne vient ni des distances,
ni seulement de l'activité individuelle non ouverte. M4 n'est jamais appelé \enquote{APL} : c'est un indice
d'accessibilité maison, meilleur que la densité.

% ---------------------------------------------------------------------------------------------------
\section{Sensibilité}\label{sec:sensibilite}
Protocole gelé et horodaté. En 2050, on compare l'effet d'une politique (quotas de formation $\times1{,}2$) à
l'étendue de $C$ entre estimateurs. Le rapport $R$ vaut :
\begin{itemize}
  \item SF 1,91\mesure{}, CD 40\mesure{}, PH 1,14\mesure{} : la politique est distinguable du choix du modèle ;
  \item MK 0,77\mesure{}, IDE 0,34\mesure{} : non distinguable.
\end{itemize}
L'étendue vient presque toute de M0 : sans lui, elle reste sous 0,004 en log\mesure{}. P-B1 ayant réfuté M0, il est
retiré de la fourchette affichée. Le verdict gelé, lui, reste publié. Les tirages de l'offre portent 52 à 92~\% de la
variance de $\log C$\mesure{}.

L'indice servi est donc M1, et M3 pour les infirmiers. L'effet du vieillissement est montré à part, et la bande
d'incertitude de l'offre sert d'incertitude.

\begin{figure}[h]\centering
\includegraphics[width=0.9\linewidth]{figures/c2050_estimateurs.pdf}
\caption{Indice de couverture national en 2050 (2026 = 1), par estimateur et par profession (grille de l'axe~5,
trajectoire centrale du jumeau, Omphale central)\mesure{}. M2a n'est pas défini pour les pharmaciens.}
\label{fig:c2050}
\end{figure}

\paragraph{Ce que change le vieillissement.} Pour les infirmiers, l'indice national 2050 vaut 1,14 si le besoin suit
la population (M0), mais 0,83 avec M1, qui pondère par l'âge\mesure{} (figure~\ref{fig:c2050}). L'offre projetée
croît plus vite que la population, mais moins vite que la demande d'une population qui vieillit. Le signe de
l'évolution s'inverse. Au département, la médiane de $C_1/C_0-1$ vaut $-28{,}4$~\% pour les infirmiers et
$-14{,}1$~\% pour les kinés\mesure{} (P-B1, annexe~\ref{ann:protocoles}). Pour les dentistes, l'âge ne change
presque rien (1,404 contre 1,405).

\begin{figure}[h]\centering
\includegraphics[width=0.9\linewidth]{figures/ratio_dep.pdf}
\caption{Rapport $R$ entre l'effet d'une hausse de 20~\% des quotas et l'étendue de $C$ entre estimateurs, par
département, en 2050 (échelle log, boîtes sans valeurs extrêmes)\mesure{}. Au-dessus de 0 : la politique se
distingue du choix du modèle.}
\label{fig:ratio}
\end{figure}

% ---------------------------------------------------------------------------------------------------
\section{Menaces à la validité}\label{sec:menaces}
\begin{itemize}
  \item \textbf{Grandeur.} On mesure la demande de ville privée, pas le besoin de la population ; la part servie par
        l'hôpital public n'est pas mesurée.
  \item \textbf{Définition de l'offre.} L'offre privée du jumeau inclut des salariés (64,5~\% pour les IDE en
        2023\mesure{}), dont les soins en clinique ne figurent pas dans le recours de ville. L'effet sur $C$ est nul si
        cette part est constante ; il ne l'est pas sur les niveaux.
  \item \textbf{Grain.} Le recours n'existe qu'à la région (DAMIR) ou en file active au département. Rien ne valide
        l'indice à la commune, sauf la comparaison avec l'APL.
  \item \textbf{Dentistes} retenus de justesse par P-B1 ; sages-femmes perdantes au backtest.
  \item \textbf{P-B6} : deux saisons avant le traitement, 12 grappes, placebo en échec.
  \item \textbf{Seuils} des protocoles fixés au gel, non étalonnés.
  \item \textbf{Paramètres non mesurables} : 14 sur 25 ; tous les transferts autres que la vaccination reposent sur
        des valeurs supposées.
\end{itemize}

% ---------------------------------------------------------------------------------------------------
\section{Reproductibilité}
\begin{itemize}
  \item Données : \texttt{v2/besoin/donnees/*.py} ; manifeste \texttt{v2/besoin/donnees/MANIFEST.md}.
  \item Estimateurs : \texttt{v2/besoin/estimateurs/} ; tests : \texttt{PYTHONPATH=v2 uv run python -m besoin.tests}
        (15/15\mesure{}).
  \item Protocoles gelés et jetons : \texttt{v2/validation/preregistered/} et \texttt{horodatage/}.
  \item Passages : \texttt{v2/besoin/validation/pb*.py}, graine 20261115.
  \item Rapports : \texttt{v2/docs/lab/reports/2026-11-15-besoin-*.md}.
  \item Dérivation : \texttt{v2/docs/lab/derivations/2026-11-15-besoin-derivation.md}.
\end{itemize}

% ---------------------------------------------------------------------------------------------------
\section{Conclusion}
En données ouvertes, le besoin de soins n'est pas identifié en niveau : la saturation n'existe pas, et la demande
induite ne se sépare pas du besoin. Ce qui reste défendable est un indice de couverture \emph{indexé}. Son
dénominateur est l'âge (M1), complété par la morbidité pour les infirmiers. Il porte l'effet du vieillissement et
rend les transferts de tâches calculables.

Les questions ouvertes ont un prix :
\begin{itemize}
  \item mesurer $k$ pour la vaccination : aucune source ouverte identifiée ;
  \item la demande servie par l'hôpital public : environ 1~h d'exploration\suppose{} ;
  \item projeter la part libérale de l'offre : environ 1~h\suppose{} ;
  \item un contrôle du distancier pour M4 : environ 45~min et un téléchargement\suppose{}.
\end{itemize}

\bibliographystyle{plainnat}
\bibliography{refs}

\clearpage
\appendix
\section{Résultats détaillés des protocoles}\label{ann:protocoles}
Toutes les valeurs sont mesurées et reprises des rapports de passage
(\texttt{v2/docs/lab/reports/2026-11-15-besoin-*.md}). Les IC sont à 95~\%, ou à 98,75~\% quand le protocole applique
la correction de Bonferroni sur quatre professions.

\subsection{P-B1 : l'âge déplace-t-il l'indice ?}
$S_1$ est la part des départements où $|C_1/C_0-1|<2$~\% en 2050. Abandon de M1 si $S_1>0{,}90$ et $\rho\ge0{,}98$,
sous P1 (poids APL) \emph{et} sous P2 (profil Open DAMIR).
\begin{center}\footnotesize
\begin{tabular}{@{}lrrrrl@{}}
\toprule
Profession & n dép. & $S_1$ & $\rho$ & médiane $C_1/C_0-1$ & Verdict \\
\midrule
Infirmiers & 100 & 0,000 & 0,835 & $-28{,}4$~\% & M1 retenu \\
Kinés & 101 & 0,000 & 0,984 & $-14{,}1$~\% & M1 retenu \\
Sages-femmes & 101 & 0,010 & 0,989 & $+6{,}1$~\% & M1 retenu \\
Dentistes & 101 & 0,931 (P1) ; 0,871 (P2) & 0,998 ; 0,999 & 0,0~\% & retenu de justesse \\
Pharmaciens & --- & 0,000 (P2) & 0,987 & --- & M1 retenu (P2 seul) \\
\bottomrule
\end{tabular}
\end{center}
Sensibilité : avec Omphale bas ou haut, $S_1$ des dentistes tombe à 0,545 et 0,584.

\subsection{P-B2 : saturation du recours}
$\beta_{Q1}$ et $\beta_{Q5}$ sont les élasticités du recours standardisé à l'accès dans le premier et le dernier
quintile, et $\delta$ leur différence. M2b tombe si l'IC de $\delta$ contient 0 ou si $\Delta$AIC $<2$.
\begin{center}\footnotesize
\begin{tabular}{@{}lrlllr@{}}
\toprule
Profession & n & $\beta_{Q1}$ [IC] & $\delta$ [IC] & $\Delta$AIC & Verdict \\
\midrule
Infirmiers & 100 & $-0{,}60$ [$-1{,}67$ ; 0,45] & 0,26 [$-1{,}66$ ; 1,63] & $-0{,}06$ & tombe \\
Kinés & 101 & 0,93 [$-0{,}25$ ; 1,97] & 0,36 [$-0{,}33$ ; 1,10] & 23,11 & tombe (non attribuable) \\
Dentistes & 101 & 1,09 [0,29 ; 1,50] & $-0{,}29$ [$-0{,}55$ ; 0,37] & 1,94 & tombe \\
Sages-femmes & 101 & 1,02 [$-0{,}12$ ; 1,59] & $-0{,}57$ [$-1{,}48$ ; 0,82] & 0,87 & tombe \\
\bottomrule
\end{tabular}
\end{center}

\subsection{P-B3 : morbidité}
$S_3$ est la part des départements où $|C_3/C_1-1|<2$~\%, sur 3 imputations des cellules masquées. Le test (B) porte
sur le $\Delta R^2$ hors échantillon de la morbidité au-delà de l'âge (seuil 0,10).
\begin{center}\footnotesize
\begin{tabular}{@{}lllll@{}}
\toprule
Profession & $S_3$ & $\rho$ & $\Delta R^2$ [IC] & Verdict \\
\midrule
Infirmiers & 0,830 & 0,992 & 0,005 [$-0{,}117$ ; 0,114] & M3 retenu pour l'indice ; H-B3 rejetée \\
Kinés & 1,000 & 0,999 & $-0{,}033$ [$-0{,}166$ ; 0,082] & recoupement ; rejetée \\
Dentistes & 0,990 & 1,000 & $-0{,}003$ [$-0{,}109$ ; 0,086] & recoupement ; rejetée \\
Sages-femmes & 0,970 & 0,999 & 0,010 [$-0{,}206$ ; 0,146] & recoupement ; rejetée \\
Pharmaciens & 0,980 & 0,999 & --- & recoupement \\
\bottomrule
\end{tabular}
\end{center}

\subsection{P-B4 et P-B4-PH : reproduction de l'APL}
EM est l'erreur relative médiane, et $G$ le gain de Spearman sur la meilleure densité naïve (commune N1 ou bassin de
vie N2). Les temps de trajet sont ceux d'OSRM.
\begin{center}\footnotesize
\begin{tabular}{@{}lllll@{}}
\toprule
Profession & $\rho_{M4}$ [IC] & EM [IC] & $G$ [IC] & Verdict \\
\midrule
Infirmiers & 0,717 [0,670 ; 0,760] & 22,7~\% [21,4 ; 24,3] & +0,147 [0,103 ; 0,204] & indice maison \\
Kinés & 0,756 [0,720 ; 0,788] & 27,5~\% [26,1 ; 29,0] & +0,191 [0,142 ; 0,252] & indice maison \\
Sages-femmes & 0,713 [0,683 ; 0,743] & 26,5~\% [24,7 ; 28,6] & +0,354 [0,318 ; 0,392] & indice maison \\
Dentistes & 0,725 [0,701 ; 0,750] & 32,8~\% [31,1 ; 34,6] & +0,299 [0,254 ; 0,345] & indice maison \\
Pharmaciens (V0 2013) & 0,723 [0,707 ; 0,738] & 33,1~\% [31,3 ; 35,3] & +0,139 [0,121 ; 0,158] & indice maison \\
\bottomrule
\end{tabular}
\end{center}
Avec le distancier à vol d'oiseau (passage 1, gelé) : $\rho$ de 0,662 à 0,766, EM de 21,2 à 29,4~\%. Le passage par la
route ne change presque rien.

\subsection{P-B5 : backtest 2016 $\to$ 2025}
MAPE départementale à 9 ans. $\Delta$ : écart de MAPE, en points, avec son IC par bootstrap des départements.
\begin{center}\footnotesize
\begin{tabular}{@{}lrrlrl@{}}
\toprule
Profession & Persistance & M1 & $\Delta_{M1}$ [IC] & GBM & $\Delta_{\mathrm{GBM}-M1}$ [IC] \\
\midrule
Infirmiers & 15,79 & 7,72 & $-8{,}07$ [$-9{,}22$ ; $-6{,}71$] & 5,81 & $-1{,}91$ [$-3{,}28$ ; $-0{,}53$] \\
Kinés & 14,11 & 10,79 & $-3{,}32$ [$-4{,}42$ ; $-2{,}13$] & 9,84 & $-0{,}95$ [$-2{,}89$ ; +0,88] \\
Dentistes & 10,37 & 9,38 & $-0{,}99$ [$-1{,}38$ ; $-0{,}56$] & 6,20 & $-3{,}19$ [$-4{,}95$ ; $-1{,}35$] \\
Sages-femmes & 65,00 & 66,57 & +1,57 [+1,35 ; +1,79] & 20,73 & $-45{,}84$ [$-51{,}34$ ; $-39{,}28$] \\
\bottomrule
\end{tabular}
\end{center}
Le gain du GBM suit la croissance de l'offre (corrélation 0,54 à 0,71 avec la croissance de l'effectif ; diagnostic a
posteriori). Il n'est servi que comme prédicteur du recours, jamais comme besoin.

\subsection{P-B6 : vaccination antigrippale en officine}
Différence de différences à adoption échelonnée : 12 zones métropolitaines, années 2014--2018, cohortes 2017
(Auvergne-Rhône-Alpes, Nouvelle-Aquitaine) et 2018 (Hauts-de-France, Occitanie). Variable :
$\log$(doses J07BB pour 1\,000 habitants).
\begin{center}\footnotesize
\begin{tabular}{@{}ll@{}}
\toprule
ATT [IC 95~\%, bootstrap sauvage] & +0,0073 [$-0{,}0012$ ; +0,0158] \\
Permutation exacte (2\,970 affectations) & $p=0{,}153$ \\
Avant traitement, $e=-3$ ; $e=-2$ & $-0{,}019$ [$-0{,}080$ ; +0,042] ; +0,004 [$-0{,}040$ ; +0,049] \\
Placebo 0--19 ans & +0,047 [+0,012 ; +0,082] : non attribuable \\
60 ans et plus ; en montants & +0,0038 [$-0{,}0029$ ; +0,0105] ; +0,0066 [$-0{,}0021$ ; +0,0153] \\
\bottomrule
\end{tabular}
\end{center}
\end{document}
